\section*{Заключение}
\addcontentsline{toc}{section}{Заключение}

В результате выполнения курсового проекта разработано веб-приложение для ведения личного 
цитатника по фильмам с системой аналитики. Подтверждено, что использование современных 
веб-технологий позволяет создать удобный инструмент для сохранения и систематизации цитат, 
а также для получения наглядной статистики по коллекции пользователя.

На основе полученных результатов можно сформулировать следующие \textbf{выводы}:

\begin{enumerate}
    \item Использование микрофреймворка Flask в сочетании с СУБД PostgreSQL и ORM SQLAlchemy 
    обеспечивает необходимую гибкость и производительность при разработке веб-приложения 
    средней сложности.
    \item Реализация агрегирующих SQL-запросов с последующей визуализацией через библиотеку 
    Chart.js позволяет наглядно отображать статистику (топ фильмов, распределение по жанрам, 
    динамику по месяцам) без излишней нагрузки на сервер.
    \item Интеграция импорта и экспорта цитат в форматах CSV и JSON даёт пользователю 
    возможность легко переносить свои данные между системами и создавать резервные копии.
    \item Разграничение прав доступа (пользователь / администратор) обеспечивает безопасность 
    данных: обычные пользователи управляют только своими цитатами, администратор поддерживает 
    общий каталог фильмов и жанров в актуальном состоянии.
\end{enumerate}

\textbf{Результаты внедрения.} Разработанное приложение развёрнуто на хостинге Render и 
доступно по адресу: \url{https://quote-app-zw4g.onrender.com/} \cite{render_docs}. Исходный код приложения размещён 
в публичном репозитории GitHub: \url{https://github.com/Kyannittee/Web_course_work_code}. Текст 
пояснительной записки (курсовой работы) в формате LaTeX также доступен в репозитории: 
\url{https://github.com/Kyannittee/Web_course_work} \cite{github_docs}. 

\textbf{Перспективы углубления темы.} В дальнейшем планируется расширить функциональность 
приложения за счёт следующих возможностей:
\begin{itemize}
    \item добавление возможности обмена цитатами между пользователями (социальная сеть);
    \item внедрение полнотекстового поиска по цитатам;
    \item реализация пагинации для оптимизации загрузки длинных списков;
    \item экспорт коллекции в PDF с кастомным оформлением.
\end{itemize}

Таким образом, все задачи, поставленные в техническом задании, выполнены в полном объёме. 
Разработанное приложение может быть использовано любителями кино для ведения личного 
цитатника, а также в качестве основы для более масштабных проектов в области сбора и 
анализа пользовательских коллекций.

При разработке пояснительной записки к курсовому проекту строго соблюдались 
требования~\cite{gost_report}, определяющие структуру и правила оформления 
научно-исследовательских работ.